In this chapter we focus on convex functions, and hence on convex optimization.

\section{General results}

We start with general results that are of use and form the basis of the theory of convex analysis.

\subsection{Definition}

\begin{definition}[Convexity for sets]
    Let $ A $ be a subset of $ \mathbf R^n $.
    
    Then, $ A $ is said to be \textit{convex} if
    \[
        \forall a \in A, \forall b \in A, \forall \lambda \in [0,1] (1-\lambda)a + \lambda b \in A
    \]
\end{definition}

\begin{remark}
    We have just defined \textit{convexity for sets}. 
    
    This notion is to be distinguished from the notion of convexity for \textit{functions}.
\end{remark}

\begin{remark}
    We will follow the convention of \cite{boydvandenberghe2004}, and often refer to the set $ \mathrm{dom}(f) $ on which $ f $ is defined (we define $ \mathrm{dom} f $ as such).
\end{remark}

\begin{definition}[Convexity for a function]
    Let $ f $ be a real valued function defined on $ \mathbf R^n $. 

    We say that $ f $ is \textit{convex} if $ \mathrm{dom} f $ is convex and if

    \[
        \forall a,b \in \mathrm{dom} f, \forall \lambda \in [0,1], f(\lambda a + (1-\lambda)b) \leqslant \lambda f(a) + (1-\lambda) f(b)
    \]
\end{definition}

\begin{definition}[Epigraph]
    The epigraph of a real-values function is

    \[ \mathrm{Epi}(f) = \{ (x,y) \in \mathrm{dom}(f) \times \mathbb R, f(x) \geqslant y \}\]
\end{definition}

\begin{proposition}
    A function is convex if and only if its epigraph is convex.
\end{proposition}

\begin{proof}
    Trivial.
\end{proof}

\subsection{The one-dimensional case}

In this section, $ f $ is defined on an open interval $ I $ of $ \mathbf R $. We have the following, which is of the uttermost importance.

\begin{theorem}[Three-slopes inequality]
    For all $ a < b < c $ in $ I $,

    \[
        \frac{f(b) - f(a)}{b-a} \leqslant \frac{f(c)-f(a)}{c-a} \leqslant \frac{f(c)-f(b)}{c-b}
    \]
\end{theorem}

% First order, <∇f(x) - ∇f(y),x-y> >= 0